\section{Encoder-Decoder and Attention}

\subsection{Why Encoder-Decoder? The Limitation of Sequence Labeling}

In sequence labeling tasks like POS tagging, each input token maps to exactly one output
token of the same position. This one-to-one alignment breaks down for tasks like machine
translation, where the source and target sentences can have different lengths and completely
different word orders. A model that pairs each input word rigidly with one output word
cannot handle this. The encoder-decoder architecture was introduced to solve this problem.

\subsection{The Three Components of an Encoder-Decoder Model}

An encoder-decoder network consists of three parts:

\begin{itemize}
    \item \textbf{Encoder:} An RNN that reads the source sequence $x_1^n$ word by word,
    updating its hidden state at each step. Its final hidden state $h_n^e$ is a compressed
    representation of the entire input sequence.

    \item \textbf{Context Vector} $c$: The final encoder hidden state passed to the decoder.
    It is the only bridge between the encoder and decoder.
    \[
        c = h_n^e
    \]

    \item \textbf{Decoder:} An RNN initialized with $c$ as its first hidden state
    $h_0^d = c$. It autoregressively generates the target sequence one token at a time,
    conditioning each step on the previous hidden state, the previous output token, and
    the context vector.
    \[
        h_t^d = g(\hat{y}_{t-1},\, h_{t-1}^d,\, c)
    \]
    \[
        z_t = f(h_t^d), \qquad y_t = \text{softmax}(z_t)
    \]
    Generation continues until an end-of-sequence marker is produced.
\end{itemize}

\subsection{The Sentence Separator Token}

To train this model, source and target sentences are concatenated with a separator token
\texttt{<s>} placed between them. The encoder processes the source side; output is ignored
during this phase. After \texttt{<s>}, the decoder begins generating the target sequence.
The model is trained to minimize cross-entropy loss averaged over all target tokens:
\[
    L = \frac{1}{T} \sum_{i=1}^{T} L_i, \qquad L_i = -\log P(y_i)
\]

\subsection{Teacher Forcing}

During inference, the decoder feeds its own predicted output $\hat{y}_t$ as input to the
next step. If an early prediction is wrong, the error compounds over subsequent steps.
During training, this is avoided using \textbf{teacher forcing}: the correct gold target
token is always used as the next input, regardless of what the decoder predicted. This
stabilizes and accelerates training.

\subsection{The Bottleneck Problem}

The context vector $c = h_n^e$ must encode the entire meaning of the source sentence into
a single fixed-size vector. For short sentences this works reasonably well, but for long
sentences information from early tokens tends to be poorly represented by the time the
encoder reaches the end. This is a \textbf{representational bottleneck} --- not a vanishing
gradient problem, but a capacity problem.

\subsection{Attention: Motivation}

Instead of relying on a single static context vector, the \textbf{attention mechanism}
lets the decoder dynamically look at all encoder hidden states at every decoding step.
Each encoder hidden state $h_j^e$ corresponds roughly to one input token. Rather than
compressing everything into one vector upfront, the decoder computes a fresh, custom
context vector $c_i$ at each step $i$, weighted toward whichever encoder states are most
relevant to the token currently being generated.

\subsection{Attention: Step-by-Step Computation}

\textbf{Step 1 --- Score each encoder state.}
At decoding step $i$, a relevance score is computed between the prior decoder hidden state
$h_{i-1}^d$ and each encoder hidden state $h_j^e$. The simplest form is
\textbf{dot-product attention}:
\[
    \text{score}(h_{i-1}^d,\, h_j^e) = h_{i-1}^d \cdot h_j^e
\]
A high score means those two vectors are similar, indicating that encoder state $j$ is
relevant to the current decoding step.

\textbf{Step 2 --- Normalize with softmax.}
Scores are converted to a probability distribution over all encoder states, giving
attention weights $\alpha_{ij}$ that sum to 1:
\[
    \alpha_{ij} = \frac{\exp(\text{score}(h_{i-1}^d,\, h_j^e))}
                       {\sum_k \exp(\text{score}(h_{i-1}^d,\, h_k^e))}
\]

\textbf{Step 3 --- Compute a weighted context vector.}
The context vector for step $i$ is a weighted average of all encoder hidden states:
\[
    c_i = \sum_j \alpha_{ij}\, h_j^e
\]

\textbf{Step 4 --- Condition the decoder on it.}
The decoder hidden state is computed using this dynamic context:
\[
    h_i^d = g(\hat{y}_{i-1},\, h_{i-1}^d,\, c_i)
\]

\subsection{Bilinear (Parameterized) Attention}

A more powerful scoring function introduces a learned weight matrix $W_s$:
\[
    \text{score}(h_{i-1}^d,\, h_j^e) = h_{i-1}^d\, W_s\, h_j^e
\]
$W_s$ is trained end-to-end alongside the rest of the model. This has two advantages
over dot-product attention: it learns \textit{which aspects} of similarity matter for the
task, and it allows the encoder and decoder to have different hidden state dimensionalities
(dot-product attention requires them to be the same size).

\subsection{Architecture of the Encoder RNN}

The encoder is not a vanilla RNN. The standard encoder design is a \textbf{stacked
bidirectional LSTM (biLSTM)}, combining three ideas:

\begin{itemize}
    \item \textbf{LSTM instead of vanilla RNN:} LSTMs use gated cell states to avoid
    vanishing gradients, allowing reliable retention of information across long sequences.
    A vanilla RNN cannot encode a long sentence reliably.

    \item \textbf{Bidirectional:} The encoder reads the source in both directions
    simultaneously. This gives every token's hidden state access to context from both
    sides of the sentence --- essential for a rich representation. The forward and backward
    hidden states are concatenated at each time step.

    \item \textbf{Stacked (multi-layer):} Multiple biLSTM layers are stacked. Lower layers
    capture surface patterns; higher layers capture abstract meaning. The context vector
    is taken from the top layer's final hidden state.
\end{itemize}

\subsection{Architecture of the Decoder RNN}

The decoder is typically a \textbf{stacked unidirectional LSTM}. It shares two of
the encoder's three properties --- LSTM gating and multi-layer stacking --- but
\textit{cannot} be bidirectional. Bidirectionality requires reading the sequence in both
directions, but the decoder generates output one token at a time during inference; future
tokens do not yet exist. The decoding direction is strictly left-to-right.

Additionally, the decoder receives inputs that a plain encoder never handles: the context
vector $c$ and the previously generated token $\hat{y}_{t-1}$ at every step.

\subsection{Why Use Two RNNs at All if the Bottleneck Problem Persists?}

The encoder-decoder bottleneck problem (compressing everything into one vector) and the
vanilla RNN gradient problem (vanishing gradients over time) are two distinct problems.
LSTMs largely solve the gradient problem. The bottleneck is a separate issue of
representational capacity, not gradient flow.

More importantly, the encoder-decoder architecture solves a problem that a single RNN
simply cannot address: generating output sequences of a different length and with
non-aligned word mappings relative to the input. A single RNN produces one output per
input token --- it has no clean mechanism for the encoder-to-decoder handoff needed for
translation. The two-RNN design made this possible. Attention was then added on top to
fix the bottleneck that remained.

The historical progression was:
\begin{enumerate}
    \item Single RNN $\to$ works for aligned same-length tasks, fails for translation.
    \item Encoder-Decoder $\to$ enables translation, struggles on long sequences.
    \item Encoder-Decoder + Attention $\to$ fixes the bottleneck, handles long sequences.
    \item Transformer $\to$ removes the RNN entirely, built purely on attention.
\end{enumerate}

\subsection{Why Bidirectionality Requires Separate Networks}

A bidirectional RNN is by definition two separate RNNs: one running forward and one
running backward, each with its own weight matrices. There is no single RNN that can
operate in both directions simultaneously --- the ``bi'' in biRNN means two networks.

Therefore, the moment you want bidirectional encoding, you have already committed to
multiple components. You cannot have a single RNN switch to bidirectional mode during
encoding and then back to unidirectional mode during decoding --- those are structurally
incompatible requirements that necessitate separate architectures.

The encoder-decoder split is not just an engineering choice. It is the necessary
architectural consequence of wanting bidirectional encoding alongside unidirectional
decoding. The two phases have fundamentally different requirements, so they require
separate components.